% TOP_PROBLEM: {"id":30007736,"problem_number":"LOCAL-30007736","title":"Returns to expanding college access","category_id":21,"category":{"id":21,"name":"economics","display_name":"Economics","description":"Open economic theory statements, published research agendas, and proposed scoped specifications; question provenance and historical priority are qualified per record.","slug":"economics","order_index":21,"created_at":"2026-10-08T00:00:00Z"},"status":"provenance_pending","external_url":null,"legacy_ids":[],"published":false,"statement":"Estimate the marginal net value of expanding one public-university system, including alternative enrollment, peer effects, instructional capacity, and existing students.","economics":{"original_id":225,"field":"Education, families and demography","kind":"empirical","formalization_basis":"proposed_specification","source_status":"not_verified","priority_status":"not_established","priority_note":"No readable primary passage explicitly stating this precise open question was verified. A supporting paper, abstract, or topic citation is insufficient to establish origin or unresolved status.","earliest_verified_reference":null,"current_reference":{"authors":["Jack Mountjoy"],"title":"Marginal Returns to Public Universities","year":2024,"url":"https://www.nber.org/papers/w32296","doi":"10.3386/w32296","locator":"Abstract","evidence_summary":"The study estimates returns on a particular admissions margin; it does not establish all margins of capacity expansion as an explicitly open problem."},"formalization_note":"Proposed scoped research specification. The original catalogue prompt was a synthesis, and no canonical conjecture or verified original formulation is claimed. The supporting source, when retained, establishes context or existing results rather than this exact open question. Mathematical scope was checked independently; added definedness, feasibility, or identification conditions belong to this proposed formalization. Mathematical scope was checked independently; added definedness, feasibility, or identification conditions belong to this proposed formalization."},"statusQualification":"Provenance pending: an explicit published statement of this exact open question was not verified. This is a catalogue-authored candidate specification, not a certified open conjecture. Proposed scoped research specification. The original catalogue prompt was a synthesis, and no canonical conjecture or verified original formulation is claimed. The supporting source, when retained, establishes context or existing results rather than this exact open question. Mathematical scope was checked independently; added definedness, feasibility, or identification conditions belong to this proposed formalization. Mathematical scope was checked independently; added definedness, feasibility, or identification conditions belong to this proposed formalization.","statusReviewedAt":"2026-10-08"}
% FIRST_POSTED: 2026-10-08
% ENTRY_KIND: statement_only
% !TeX program = lualatex
\documentclass[11pt]{article}
\usepackage{fontspec}
\usepackage[a4paper,margin=25mm]{geometry}
\usepackage{amsmath,amssymb,mathtools}
\usepackage{xurl}
\usepackage[hidelinks]{hyperref}
\newcommand{\sref}[2]{\hyperlink{source:#1:#2}{[S#2]}}
\setlength{\emergencystretch}{3em}
\begin{document}
% problemId:30007736
\section*{Returns to expanding college access}\label{30007736}
\subsection{Definitions and mathematical statement}
Consider applicants near the admission threshold of one public-university system. Let \(k\in\{0,1,\ldots,\bar k\}\) be an administratively feasible integer capacity expansion, where \(\bar k\) is a declared maximum measured in additional seats, and let \(E_i(k)\) be applicant \(i\)'s enrollment destination after admissions, financial aid, and peer composition adjust. Define discounted lifetime earnings \(Y_i(k)\), education resource cost \(C_i(k)\), and a prespecified monetary nonpecuniary benefit \(N_i(k)\). For the baseline applicant cohort, define
\[V(k)=E\!\left[\sum_i\{Y_i(k)+N_i(k)-C_i(k)\}\right]-K(k),\]
where \(K(k)\) is additional infrastructure cost not already included in \(C_i(k)\). Determine the discrete marginal net value \(V(k+1)-V(k)\) for \(k<\bar k\) and distinguish it from a local admissions-cutoff earnings effect. The identification strategy must address alternative colleges, displacement of other applicants, completion, peer effects, and declining instructional resources per pupil. Lifetime quantities beyond observed follow-up require explicit extrapolation and uncertainty intervals. Report separately outcomes for marginal admits and changes for existing students. The cited study estimates a specific university-admissions margin; it does not by itself establish this complete capacity-expansion counterfactual as an explicit unresolved question. The proposed welfare target therefore has no verified priority claim.

\par\textbf{Formulation and scope.} Proposed scoped research specification. The original catalogue prompt was a synthesis, and no canonical conjecture or verified original formulation is claimed. The supporting source, when retained, establishes context or existing results rather than this exact open question. Mathematical scope was checked independently; added definedness, feasibility, or identification conditions belong to this proposed formalization. Mathematical scope was checked independently; added definedness, feasibility, or identification conditions belong to this proposed formalization.

\par\textbf{Open-question provenance.} An explicit published open-question statement was not verified. Sources below are supporting literature and must not be treated as a first listing.

\par\textbf{Historical priority.} No readable primary passage explicitly stating this precise open question was verified. A supporting paper, abstract, or topic citation is insufficient to establish origin or unresolved status.

\subsection{Short English statement}
Estimate the marginal net value of expanding one public-university system, including alternative enrollment, peer effects, instructional capacity, and existing students.

\subsection{Sources}
\begin{itemize}\raggedright
\item[\textnormal{[S1]}]\hypertarget{source:30007736:1}{} Jack Mountjoy. 2024. Marginal Returns to Public Universities. Abstract.
\url{https://www.nber.org/papers/w32296}.
DOI: 10.3386/w32296.
{\small\textit{Supporting or current literature; see evidence qualification. The study estimates returns on a particular admissions margin; it does not establish all margins of capacity expansion as an explicitly open problem.}}
\end{itemize}
\end{document}
